% Part: second-order-logic
% Chapter: syntax-and-semantics
% Section: language-of-sol

\documentclass[../../../include/open-logic-section]{subfiles}

\begin{document}

\olfileid[hi]{sol}{syn}{lan}

\olsection{द्वितीय-क्रम तर्कशास्त्र की भाषा}

प्रथम-क्रम तर्कशास्त्र की तरह द्वितीय-क्रम तर्कशास्त्र की अभिव्यक्तियाँ एक
मूल शब्दावली से बनती हैं जिसमें \emph{!!{variable}s},
\emph{!!{constant}s}, \emph{!!{predicate}s} और कभी-कभी
\emph{!!{function}s} होते हैं। इनसे, तार्किक संयोजकों, परिमाणकों तथा कोष्ठक
और अल्पविराम जैसे विराम-चिह्नों के साथ, \emph{पद} और \emph{!!{formula}s}
बनते हैं। अंतर यह है कि वस्तुओं के चरों के अतिरिक्त द्वितीय-क्रम तर्कशास्त्र
में संबंधों और फलनों के चर भी होते हैं तथा उन पर परिमाणीकरण की अनुमति होती
है। अतः द्वितीय-क्रम तर्कशास्त्र के तार्किक चिह्न प्रथम-क्रम तर्कशास्त्र के
चिह्नों के साथ ये हैं:

\begin{enumerate}
\item प्रत्येक स्थान-संख्या~$n$ के द्वितीय-क्रम संबंध !!{variable}s का एक
  !!{denumerable} समुच्चय: $\Obj V_0^n$, $\Obj V_1^n$, $\Obj V_2^n$, \dots
\item द्वितीय-क्रम फलन !!{variable}s का एक !!{denumerable} समुच्चय:
  $\Obj u_0^n$, $\Obj u_1^n$, $\Obj u_2^n$, \dots
\end{enumerate}

प्रथम-क्रम तर्कशास्त्र में प्रायः !!{identity}~$\eq$ सम्मिलित होती है।
प्रथम-क्रम तर्कशास्त्र में किसी भाषा~$\Lang{L}$ के अतार्किक चिह्न हमें कोई
रोचक बात व्यक्त करने देने के लिए निर्णायक हैं। निस्संदेह ऐसे !!{sentence}s
हैं जिनमें कोई अतार्किक चिह्न नहीं आता, किंतु केवल~$\eq$ से कुछ रोचक कहना
कठिन है। द्वितीय-क्रम तर्कशास्त्र में संबंध और फलन चरों की असीमित आपूर्ति होने
के कारण हम अतार्किक चिह्नों की विशेष आपूर्ति के बिना भी वह सब कह सकते हैं जो
किसी प्रथम-क्रम भाषा में कह सकते हैं।

\end{document}
